\oldpage[I]{85}
Lorsque $f$ parcourt $A-\mathfrak{j}_x$ (pour un $x$ donné dans
$X=\operatorname{Spec}(A)$), les ensembles $S'_f$ constituent un ensemble
filtrant croissant de parties de $A-\mathfrak{j}_x$, car pour deux éléments
$f$, $g$ de $A-\mathfrak{j}_x$, $S'_f$ et $S'_g$ sont contenus dans
$S'_{fg}$~; comme la réunion des $S'_f$ pour $f\in A-\mathfrak{j}_x$ est
$A-\mathfrak{j}_x$, on en conclut (\hyperref[0.1.4.5-fr]{0, 1.4.5}) que le
$A_x$-module $M_x$ s'identifie canoniquement à la \emph{limite inductive}
$\varinjlim M_f$, relativement à la famille d'homomorphismes
$(\rho_{g,f})$. Nous désignerons par
\[
  \rho_x^f:M_f\longrightarrow M_x
\]
l'homomorphisme canonique pour $f\in A-\mathfrak{j}_x$ (ou, ce qui revient au
même, $x\in D(f)$).

\begin{definition}[1.3.4]
\label{I.1.3.4-fr}
On appelle \emph{faisceau structural du spectre premier}
$X=\operatorname{Spec}(A)$ (resp. \emph{faisceau associé au $A$-module}
$M$) et on note $\widetilde{A}$ ou $\mathscr{O}_X$ (resp. $\widetilde{M}$) le
faisceau d'anneaux (resp. le $\widetilde{A}$-Module) associé au préfaisceau
$D(f)\mapsto A_f$ (resp. $D(f)\mapsto M_f$) sur la base $\mathfrak{B}$ de
$X$ formée des $D(f)$ où $f\in A$
((\hyperref[I.1.1.10-fr]{1.1.10}),
(\hyperref[0.3.2.1-fr]{0, 3.2.1} et
\hyperref[0.3.5.6-fr]{3.5.6})).
\end{definition}

On a vu (\hyperref[0.3.2.4-fr]{0, 3.2.4}) que la fibre
$\widetilde{A}_x$ (resp. $\widetilde{M}_x$) \emph{s'identifie à l'anneau}
$A_x$ (resp. \emph{au $A_x$-module} $M_x$)~; nous désignerons par
\[
  \theta_f:A_f\longrightarrow\Gamma(D(f),\widetilde{A})
\]
\[
  \text{(resp. }\theta_f:M_f\longrightarrow
  \Gamma(D(f),\widetilde{M})\text{),}
\]
l'application canonique, de sorte que pour tout $x\in D(f)$ et tout
$\xi\in M_f$, on a
\begin{equation}
  (\theta_f(\xi))_x=\rho_x^f(\xi).
  \tag{1.3.4.1}\label{I.1.3.4.1-fr}
\end{equation}

\begin{proposition}[1.3.5]
\label{I.1.3.5-fr}
$\widetilde{M}$ est un foncteur covariant exact en $M$, de la catégorie des
$A$-modules dans la catégorie des $\widetilde{A}$-Modules.
\end{proposition}

En effet, soient $M$, $N$ deux $A$-modules, $u$ un homomorphisme
$M\to N$~; pour tout $f\in A$, il correspond canoniquement à $u$ un
homomorphisme $u_f$ du $A_f$-module $M_f$ dans le $A_f$-module $N_f$, et le
diagramme (pour $D(g)\subset D(f)$)
\[
  \xymatrix{
    M_f\ar[r]^{u_f}\ar[d]_{\rho_{g,f}} & N_f\ar[d]^{\rho_{g,f}}\\
    M_g\ar[r]_{u_g} & N_g
  }
\]
est commutatif (\hyperref[0.1.4.1-fr]{0, 1.4.1})~; ces homomorphismes
définissent donc un homomorphisme de $\widetilde{A}$-Modules
$\widetilde{u}:\widetilde{M}\to\widetilde{N}$
(\hyperref[0.3.2.3-fr]{0, 3.2.3}). En outre, pour tout $x\in X$,
$\widetilde{u}_x$ est la limite inductive des $u_f$ pour $x\in D(f)$
($f\in A$), et par suite (\hyperref[0.1.4.5-fr]{0, 1.4.5}) si on identifie
canoniquement $\widetilde{M}_x$ et $\widetilde{N}_x$ à $M_x$ et $N_x$
respectivement, $\widetilde{u}_x$ s'identifie à l'homomorphisme $u_x$ déduit
canoniquement de $u$. Si $P$ est un troisième $A$-module, $v$ un
homomorphisme $N\to P$ et $w=v\circ u$, il est immédiat que
$w_x=v_x\circ u_x$, donc $\widetilde{w}=\widetilde{v}\circ\widetilde{u}$. On
a donc bien défini un \emph{foncteur covariant} $\widetilde{M}$ en $M$, de la
catégorie des $A$-modules dans celle des $\widetilde{A}$-Modules. \emph{Ce
foncteur est exact}, car pour tout $x\in X$, $M_x$ est un foncteur exact en
$M$ (\hyperref[0.1.3.2-fr]{0, 1.3.2})~; en outre, on a
$\operatorname{Supp}(M)=\operatorname{Supp}(\widetilde{M})$ en vertu des
définitions des deux membres (\hyperref[0.1.7.1-fr]{0, 1.7.1} et
\hyperref[0.3.1.6-fr]{3.1.6}).

